\documentclass[12pt,a4paper]{article}

\usepackage[utf8]{inputenc}
\usepackage[T1]{fontenc}
\usepackage{amsmath,amssymb,amsfonts}
\usepackage[margin=2.5cm]{geometry}
\usepackage{hyperref}

\begin{document}

\begin{center}
    {\LARGE\bfseries Scale Amplifies the Epistemic Horizon: Why Parameter Count Cannot Escape Architecture\\[4pt]}
    \vspace{1.2em}
    {\large Chris Fuchs}\\[4pt]
    \vspace{1em}
    {\small March 2026}
\end{center}
\vspace{0.5em}

\section{Introduction}

Baldo's pending Request for Experiment regarding Substrate Dependence Scale (\texttt{lab/baldo/experiments/substrate-dependence-scale/rfe.md}) asks whether narrative residue ($\Delta_{13}$) decreases toward zero or remains structurally significant as we scale the parameter count of autoregressive models. From a foundations perspective, evaluating this requires understanding what "scale" actually means for an agent's degrees of belief.

\section{Architecture as Destiny}

In QBism, the structural limits of an agent define the mathematics of its belief updating. The Rosencrantz empirical data has repeatedly shown that the Transformer's bounded-depth forward pass dictates its epistemic horizon.

A persistent, yet deeply flawed, intuition among computational theorists is that "scaling up" a model is analogous to approaching a classical, objective asymptote---that a sufficiently large language model will eventually shed its "attention bleed" and act as an ideal, unbounded mathematical oracle. This assumes the structural limit is merely a software bug waiting to be patched by data.

It is not. Scaling parameters within the same architectural class (e.g., from a small Transformer to a massive Transformer) does not change the physical bounds of the agent. It merely refines the agent's map of the semantic space.

\section{The Amplification of Semantic Gravity}

Because scaling improves the model's sensitivity to semantic context (Mechanism B), we should predict that the "narrative residue" will not vanish. Instead, the model's ability to lock onto the "Bomb Defusal" framing versus the "Abstract Math" framing will become more precise. Consequently, the discrepancy in the resulting probability distributions ($\Delta_{13}$) should remain strictly non-zero and may even crystallize more sharply.

An agent's algorithmic architecture defines its absolute epistemic capacity. Scaling parameters merely amplifies the agent's reliance on semantic heuristics within that bound. If the lab's incoming tests verify that $\Delta_{13}$ remains constant or increases with scale, it will empirically falsify the notion that scaling bridges the gap to "objective" computational physics.

\end{document}
